% GENERATED FILE. Edit canonical lessons, then run npm run generate:books.
% canonical-source-hash: fnv1a64:a5bb8c0f5e92599f
% canonical-lessons: PT-C503-xadrez, PT-C503-baralho, PT-C503-tabuleiro, PT-C503-puzzle, PT-C503-palavras-cruzadas

\chapter{Jogos de mesa}
\label{ch:pt-jogos-de-mesa}
\hlchaptermodality{503}

\begin{chapteropening}
\textbf{By the end of this chapter:} \emph{I can talk about board games, cards and puzzles.}

Complete the last lesson of chapter 503: I can talk about board games, cards and puzzles.
\end{chapteropening}

\section[o xadrez]{o xadrez — chess; check pattern}

\label{lesson:PT-C503-xadrez}



\begin{quote}
Before the new one: say the Portuguese for fish hook, then the Portuguese for to hunt.
\end{quote}

\subsection*{You'll want to know: o xadrez}
\textbf{o xadrez} — \textquotedblleft{}chess; check pattern\textquotedblright{}.

\begin{quote}
O meu avô ensinou-me a jogar xadrez.
\end{quote}

\begin{quote}
\emph{My grandfather taught me to play chess.}
\end{quote}

Also a check pattern: \emph{uma camisa de xadrez}.

\begin{cousinweb}
Like its Spanish and Mirandese cousins, \emph{xadrez} came through Andalusian Arabic from Arabic \emph{šaṭranj} \textquotedblleft{}chess\textquotedblright{}, which goes back to Persian (EtymDB); the Persian word is usually traced to Sanskrit \emph{caturaṅga} \textquotedblleft{}four-limbed\textquotedblright{}, the name of an army.
\end{cousinweb}

\begin{grammarlens}[title={What you've built}]
One more word for telling, explaining and describing at B1.
\end{grammarlens}

\subsection*{Your turn}
\begin{itemize}
\raggedright
  \item \emph{Say it:} \emph{o xadrez}
  \item \emph{Say it:} \emph{o xadrez}, once more
  \item \emph{Say it:} a sentence of your own with \emph{o xadrez}
  \item \emph{Recall:} say the Portuguese for fish hook, then the Portuguese for to hunt, then say \emph{o xadrez} again
\end{itemize}

\begin{tcolorbox}[breakable,colback=teal!4,colframe=teal!35!black,title={Before you move on}]
What does o xadrez mean? (\textquotedblleft{}Chess; check pattern\textquotedblright{}.) Say it once more.
\end{tcolorbox}
\section[o baralho]{o baralho — pack of cards, deck}

\label{lesson:PT-C503-baralho}



\begin{quote}
Before the new one: say the Portuguese for to hunt, then the Portuguese for chess.
\end{quote}

\subsection*{You'll want to know: o baralho}
\textbf{o baralho} — \textquotedblleft{}pack of cards, deck\textquotedblright{}.

\begin{quote}
Temos um baralho de cartas na gaveta.
\end{quote}

\begin{quote}
\emph{We have a pack of cards in the drawer.}
\end{quote}

\emph{Baralhar} — to shuffle; also to muddle.

\begin{cousinweb}
The origin of \emph{baralho} is not settled; it may be linked to \emph{baralhar} \textquotedblleft{}to mix up\textquotedblright{}, but the source of both is uncertain.
\end{cousinweb}

\begin{grammarlens}[title={What you've built}]
One more word for telling, explaining and describing at B1.
\end{grammarlens}

\subsection*{Your turn}
\begin{itemize}
\raggedright
  \item \emph{Say it:} \emph{o baralho}
  \item \emph{Say it:} \emph{o baralho}, once more
  \item \emph{Say it:} a sentence of your own with \emph{o baralho}
  \item \emph{Recall:} say the Portuguese for to hunt, then the Portuguese for chess, then say \emph{o baralho} again
\end{itemize}

\begin{tcolorbox}[breakable,colback=teal!4,colframe=teal!35!black,title={Before you move on}]
What does o baralho mean? (\textquotedblleft{}Pack of cards, deck\textquotedblright{}.) Say it once more.
\end{tcolorbox}
\section[o tabuleiro]{o tabuleiro — (game) board; tray}

\label{lesson:PT-C503-tabuleiro}



\begin{quote}
Before the new one: say the Portuguese for chess, then the Portuguese for pack of cards.
\end{quote}

\subsection*{You'll want to know: o tabuleiro}
\textbf{o tabuleiro} — \textquotedblleft{}(game) board; tray\textquotedblright{}.

\begin{quote}
As peças do xadrez estão no tabuleiro.
\end{quote}

\begin{quote}
\emph{The chess pieces are on the board.}
\end{quote}

Also a tray: \emph{um tabuleiro de bolos}.

\begin{cousinweb}
\emph{tabuleiro} is inherited, through Old Portuguese \emph{tavoleiro} \textquotedblleft{}gameboard\textquotedblright{} (EtymDB), from a form built on Latin \emph{tabula} \textquotedblleft{}a board, a plank\textquotedblright{}.
\end{cousinweb}

\begin{grammarlens}[title={What you've built}]
One more word for telling, explaining and describing at B1.
\end{grammarlens}

\subsection*{Your turn}
\begin{itemize}
\raggedright
  \item \emph{Say it:} \emph{o tabuleiro}
  \item \emph{Say it:} \emph{o tabuleiro}, once more
  \item \emph{Say it:} a sentence of your own with \emph{o tabuleiro}
  \item \emph{Recall:} say the Portuguese for chess, then the Portuguese for pack of cards, then say \emph{o tabuleiro} again
\end{itemize}

\begin{tcolorbox}[breakable,colback=teal!4,colframe=teal!35!black,title={Before you move on}]
What does o tabuleiro mean? (\textquotedblleft{}(game) board; tray\textquotedblright{}.) Say it once more.
\end{tcolorbox}
\section[o puzzle]{o puzzle — jigsaw puzzle}

\label{lesson:PT-C503-puzzle}



\begin{quote}
Before the new one: say the Portuguese for pack of cards, then the Portuguese for (game) board.
\end{quote}

\subsection*{You'll want to know: o puzzle}
\textbf{o puzzle} — \textquotedblleft{}jigsaw puzzle\textquotedblright{}.

\begin{quote}
Fizemos um puzzle de mil peças nas férias.
\end{quote}

\begin{quote}
\emph{We did a thousand-piece jigsaw on holiday.}
\end{quote}

In Brazil: \emph{o quebra-cabeça}.

\begin{cousinweb}
\emph{puzzle} was borrowed from English \emph{puzzle}; the English word has no settled origin.
\end{cousinweb}

\begin{grammarlens}[title={What you've built}]
One more word for telling, explaining and describing at B1.
\end{grammarlens}

\subsection*{Your turn}
\begin{itemize}
\raggedright
  \item \emph{Say it:} \emph{o puzzle}
  \item \emph{Say it:} \emph{o puzzle}, once more
  \item \emph{Say it:} a sentence of your own with \emph{o puzzle}
  \item \emph{Recall:} say the Portuguese for pack of cards, then the Portuguese for (game) board, then say \emph{o puzzle} again
\end{itemize}

\begin{tcolorbox}[breakable,colback=teal!4,colframe=teal!35!black,title={Before you move on}]
What does o puzzle mean? (\textquotedblleft{}Jigsaw puzzle\textquotedblright{}.) Say it once more.
\end{tcolorbox}
\section[as palavras cruzadas]{as palavras cruzadas — crossword}

\label{lesson:PT-C503-palavras-cruzadas}



\begin{quote}
Before the new one: say the Portuguese for (game) board, then the Portuguese for jigsaw puzzle.
\end{quote}

\subsection*{You'll want to know: as palavras cruzadas}
\textbf{as palavras cruzadas} — \textquotedblleft{}crossword\textquotedblright{}.

\begin{quote}
O meu pai faz as palavras cruzadas do jornal todos os dias.
\end{quote}

\begin{quote}
\emph{My father does the newspaper's crossword every day.}
\end{quote}

Literally \textquotedblleft{}crossed words\textquotedblright{}. Plural in Portuguese.

\begin{cousinweb}
\emph{palavra} is inherited from Latin \emph{parabola} \textquotedblleft{}a comparison, a word\textquotedblright{}, and \emph{cruzada} is built on \emph{cruz} \textquotedblleft{}cross\textquotedblright{}, inherited from Latin \emph{crux} (Lewis \& Short).
\end{cousinweb}

\begin{grammarlens}[title={What you've built}]
One more word for telling, explaining and describing at B1.
\end{grammarlens}

\subsection*{Your turn}
\begin{itemize}
\raggedright
  \item \emph{Say it:} \emph{as palavras cruzadas}
  \item \emph{Say it:} \emph{as palavras cruzadas}, once more
  \item \emph{Say it:} a sentence of your own with \emph{as palavras cruzadas}
  \item \emph{Recall:} say the Portuguese for (game) board, then the Portuguese for jigsaw puzzle, then say \emph{as palavras cruzadas} again
\end{itemize}

\begin{tcolorbox}[breakable,colback=teal!4,colframe=teal!35!black,title={Before you move on}]
What does as palavras cruzadas mean? (\textquotedblleft{}Crossword\textquotedblright{}.) Say it once more.
\end{tcolorbox}
